% sections/04-preguntas-adicionales.tex
\section{Preguntas adicionales}

\pregunta{¿Qué es una inyección SQL?}
\begin{respuesta}
\end{respuesta}

\porque{base para la respuesta}{%
  ``Es una vulnerabilidad de la capa de aplicación en la que una entrada no
  confiable del usuario se concatena directamente sobre una sentencia SQL sin
  sanitizar ni parametrizar, de forma que el atacante logra alterar la
  semántica de la consulta que el motor termina ejecutando, pudiendo eludir
  esquemas de autenticación, exponer o sobrescribir información de la base de
  datos y, en escenarios con privilegios suficientes, incluso llegar a
  comandos del sistema operativo del anfitrión.'' Apóyate en el marco
  conceptual de la sección 2 y cita \texttt{owasp-sqli}.%
}

\pregunta{¿Qué es una inyección SQL activa y pasiva? Establezca las diferencias y los métodos como se ejecutan.}
\begin{respuesta}
\end{respuesta}

\porque{base para la respuesta}{%
  En la \textbf{activa} el atacante envía sentencias y observa directamente la
  respuesta del servidor (datos visibles en pantalla o mensajes de error),
  típicamente apoyado en herramientas automatizadas de fuzzing como sqlmap,
  que prueban diccionarios de payloads contra cada parámetro, tal como se hizo
  en la sección 3; en la \textbf{pasiva} el atacante no altera respuestas sino
  que inyecta sobre puntos de entrada ya existentes (como un formulario de
  login) y deduce el efecto del comportamiento de la aplicación, sin obtener
  necesariamente un volcado directo de datos, tal como se hizo en la sección 4.
  Diferencias a listar: grado de interacción con la respuesta, uso o no de
  herramientas, ruido generado (la activa es mucho más detectable en logs) y
  método (payloads automatizados y observación directa vs.\ inyección manual e
  inferencia). Cuidado: en la literatura hay varias taxonomías (blind, error-
  based, etc.); si el profesor manejó otra clasificación en clase, alinea la
  respuesta con esa.%
}

\pregunta{¿Cómo es el funcionamiento de la herramienta sqlmap?}
\begin{respuesta}
\end{respuesta}

\porque{base para la respuesta}{%
  ``sqlmap automatiza el proceso completo de la inyección: primeramente hace
  fingerprinting del objetivo (servidor web y DBMS, con pruebas específicas
  por motor), posteriormente prueba diccionarios de payloads contra cada
  parámetro de la URL o de los formularios, identificando la técnica que
  mejor responde (union query, error-based, boolean-based, time-based), y con
  el canal de inyección establecido va encadenando sus opciones de
  post-explotación (\texttt{--banner}, \texttt{--passwords}, \texttt{--dump},
  entre otras) para enumerar y extraer la información del gestor.'' Cita
  \texttt{sqlmap-docs} y relationa cada fase con lo evidenciado en las
  capturas de la sección 3.%
}

\pregunta{¿Qué controles puede implementar para evitar la inyección de sentencias SQL sobre aplicaciones web?}
\begin{respuesta}
\end{respuesta}

\porque{base para la respuesta}{%
  Lista con término en negrita + explicación, tu patrón usual:
  \textbf{Sentencias preparadas y consultas parametrizadas:} tratar la entrada
  siempre como dato y nunca como parte del código SQL, que es el control
  fundamental; \textbf{Validación de entrada allowlist:} rechazar caracteres o
  formatos que no correspondan al tipo esperado del campo; \textbf{Mínimo
  privilegio:} que el usuario de conexión de la aplicación solo acceda al
  esquema que necesita, sin permisos de administración ni de lectura sobre
  tablas ajenas; \textbf{Gestión de errores:} no exponer al usuario mensajes
  de error del motor que revelen estructura interna; \textbf{WAF:} filtrado
  de payloads conocidos como capa complementaria (nunca sustituto de la
  parametrización); \textbf{Actualización y parches:} del motor y frameworks
  para cerrar vectores conocidos. Cierra citando \texttt{owasp-sqli}.%
}
